\documentclass[UTF8,zihao=-4]{ctexart}
\usepackage[a4paper,margin=2.5cm]{geometry}
\usepackage{booktabs,tabularx,array}
\usepackage{hyperref}
\hypersetup{hidelinks}
\setlength{\parindent}{2em}
\linespread{1.25}
\renewcommand{\arraystretch}{1.2}
\pagestyle{plain}

\begin{document}
\begin{center}
{\zihao{2}\bfseries AI工具使用详情\par}
\vspace{0.8em}
{\zihao{-4}全国大学生数学建模竞赛支撑材料}
\end{center}

\section{使用的工具}

本队使用OpenAI Codex辅助整理本地项目中的题面、建模文档、程序结果和论文格式要求。使用日期为2026年9月；任务界面显示模型为GPT-5系列。AI仅作为辅助工具，核心模型选择、参数口径、程序运行、结果判断和最终提交责任均由参赛队承担。

\section{使用目的和环节}

{\small
\setlength{\tabcolsep}{4pt}
\begin{tabular}{p{1.7cm}p{6.0cm}p{5.2cm}}
\toprule
环节 & AI辅助内容 & 人工责任 \\
\midrule
论文结构 & 根据2026官方格式建立无目录的摘要、正文、声明、参考文献和附录结构 & 队员核对赛区要求和最终页数 \\
语言与排版 & 将已有模型与结果整理为中文科技论文表述，生成LaTeX排版源文件 & 队员逐句确认含义、删改表述 \\
模型审阅 & 检查功率/电量、信息时间线、跨日SOC和结算变量是否一致 & 建模手确认公式与最终实现相符 \\
结果整理 & 从正式结果说明和结果文件提取四问数值与图表，不生成缺失数值 & 编程手独立重算费用和约束 \\
合规检查 & 提示AI声明、支撑材料、匿名性和附件要求 & 全队完成最终提交检查 \\
\bottomrule
\end{tabular}
}

\section{主要提示方式与过程}

主要提示要求AI在指定项目目录内工作，先读取题面分析、四问建模文档、程序输出、复核记录和论文格式文件，再根据官方规范完善论文工程；要求所有数值只能来自团队已有正式计算记录，不能虚构结果；要求严格区分决策时可用信息与事后实际数据，并保留统计不确定性和不利结果。AI随后核对2026年官方论文格式与AI使用规定，按提交时间辨认最新成果，整理LaTeX正文、图表、附录源程序与本说明。

\section{采纳、修改和核验}

论文采纳了AI辅助生成的章节组织、部分语言表述和LaTeX排版代码。四问的数值、图表和测试记录均来自团队已有正式成果；AI没有替代程序求解，也没有把未通过事前验证的事后较优参数改成主方案。对于AI建议，团队采用如下人工复核流程：

\begin{enumerate}
  \item 建模手逐式核对公共约束、四问目标函数与信息时间线，特别检查未来信息泄漏。
  \item 编程手以固定配置复核运行记录，对费用、能量平衡、SOC和结果工作簿作独立重算。
  \item 论文手逐句审查摘要、结论、引用和所有百分比的比较基准，删除无法由结果支撑的表述。
  \item 全队在提交前再次检查匿名性、正文页数、PDF大小、源程序完整性和支撑材料文件树。
\end{enumerate}

AI输出可能存在文字误解、数值抄录和格式错误，团队已要求每项结论可追溯到正式结果文件；本说明不免除参赛队对最终作品真实性、原创性和合规性的责任。

\end{document}
